% chapters/10/checklist.tex
\bagian{Checklist Pencapaian Kompetensi}

Bab ini menuntaskan \SUBCPMK{2.3}. Centang setiap butir berikut hanya apabila
Anda benar-benar telah melakukannya atau dapat menjelaskannya.

\begin{checklistbox}
\begin{checklist}
  \item {\bfseries (Indikator a)} Saya dapat menjelaskan bagaimana peta
        berpetak diubah menjadi graf, beserta biaya setiap sisinya.
  \item Saya dapat menulis fungsi penentuan tetangga yang menyaring petak di
        luar batas peta dan petak dinding.
  \item Program saya memeriksa petanya sebelum dipakai, yaitu panjang baris
        serta keberadaan titik awal dan titik tujuan.
  \item {\bfseries (Indikator b)} Saya dapat menjelaskan perbedaan BFS,
        Dijkstra, dan A* beserta keadaan yang membuat masing-masing diperlukan.
  \item Saya dapat menjelaskan arti $g$, $h$, dan $f = g + h$, serta alasan
        urutan pemeriksaan ditentukan oleh $f$.
  \item Saya dapat menjelaskan syarat yang harus dipenuhi sebuah heuristik agar
        jalur yang dihasilkan tetap terpendek.
  \item Saya telah mengukur ketiga heuristik pada peta yang sama dan
        menunjukkan bahwa biaya jalurnya sama sedangkan banyaknya petak yang
        diperiksa berbeda.
  \item Saya dapat menjelaskan mengapa jarak Manhattan lebih mengarahkan
        pencarian daripada jarak Euclidean pada gerak empat arah.
  \item Saya dapat menjelaskan akibat heuristik yang melebihkan jarak terhadap
        mutu jalur, dan bagaimana mengenalinya.
  \item {\bfseries (Indikator c)} Saya dapat menulis rekonstruksi jalur dari
        tabel asal, beserta penanganan keadaan tujuan yang tidak dapat dicapai.
  \item Saya dapat menulis pemeriksaan keabsahan jalur, yaitu bersambung,
        dimulai dari titik awal, berakhir di titik tujuan, dan tidak melewati
        dinding.
  \item Program saya memvisualisasikan petak yang telah diperiksa, petak yang
        menunggu diperiksa, dan jalur yang ditemukan.
  \item Saya dapat menyimpan tangkapan layar keadaan pencarian dan
        memakainya pada laporan praktikum.
  \item Catatan praktikum saya memuat analisis pengaruh heuristik terhadap
        kinerja pencarian.
\end{checklist}
\end{checklistbox}
